\section{Binomial distribution}
\label{binomialModel}

\index{distribution!binomial|(}

The \termsub{binomial distribution}{distribution!binomial}
is used to describe
the number of successes in a fixed number of trials.
%,
%and this distribution is occasionally used in statistics,
%especially when doing more careful analysis of samples
%of data where simpler tools are not helpful.
This is different from the geometric distribution,
which described the number of trials we must wait before
we observe a success.


\subsection{The binomial distribution}

%\newcommand{\insureSprob}{0.7}
%\newcommand{\insureSperc}{70\%}
%\newcommand{\insureFprob}{0.3}
%\newcommand{\insureFperc}{30\%}
%\newcommand{\insureDistA}{0.7}
%\newcommand{\insureDistB}{0.21}
%\newcommand{\insureDistC}{0.063}
%\newcommand{\insureDistD}{0.019}
%\newcommand{\insureDistE}{0.006}
%\newcommand{\insureCDistA}{0.7}
%\newcommand{\insureCDistB}{0.91}
%\newcommand{\insureCDistC}{0.973}
%\newcommand{\insureCDistCComplement}{0.027}
%\newcommand{\insureCDistD}{0.992}
%\newcommand{\insureCDistE}{0.998}
%\newcommand{\insureGeomMean}{1.43}
\newcommand{\insureS}{\resp{not}}
\newcommand{\insureF}{\resp{exceed}}
% Doesn't consider binomial coefficient in next calculated value.
\newcommand{\insureBinomCinDSingleScenario}{0.103}
\newcommand{\insureBinomCinD}{0.412}
\newcommand{\insureBinomEinHSingleScenario}{0.00454}
\newcommand{\insureBinomEinH}{0.254}
\newcommand{\insureBinomFourtyExpValue}{28}
\newcommand{\insureBinomFourtySD}{2.9}
\newcommand{\insureBinomFourtyLower}{22}
\newcommand{\insureBinomFourtyUpper}{34}

\noindent%
Let's again imagine ourselves back at the insurance agency
where \insureSperc{} of individuals do not exceed their
deductible.

\begin{examplewrap}
\begin{nexample}{Suppose the insurance agency is considering
    a random sample of four individuals they insure.
    What is the chance exactly one of them will exceed
    the deductible and the other three will not?
    Let's call the four people
    Ariana ($A$),
    Brittany ($B$),
    Carlton ($C$),
    and Damian ($D$)
    for convenience.}
  \label{insureOneOfFourExceedsDeductible}%
  Let's consider a scenario where one person exceeds
  the deductible:
  \begin{align*}
  &P(A=\text{\insureF{}},
      \text{ }B=\text{\insureS{}},
      \text{ }C=\text{\insureS{}},
      \text{ }D=\text{\insureS{}}) \\
    &\quad = P(A=\text{\insureF{}})\ 
        P(B=\text{\insureS{}})\ 
        P(C=\text{\insureS{}})\ 
        P(D=\text{\insureS{}}) \\
    &\quad =  (\insureFprob{})
        (\insureSprob{})
        (\insureSprob{})
        (\insureSprob{}) \\
    &\quad = (\insureSprob{})^3 (\insureFprob{})^1 \\
    &\quad = \insureBinomCinDSingleScenario{}
  \end{align*}
  But there are three other scenarios: Brittany, Carlton,
  or Damian could have been the one to exceed the deductible.
  In each of these cases, the probability is again
  $(\insureSprob{})^3 (\insureFprob{})^1$.
  These four scenarios exhaust all the possible ways that
  exactly one of these four people could have exceeded
  the deductible, so the total probability is
  $4 \times (\insureSprob{})^3 (\insureFprob{})^1
      = \insureBinomCinD{}$.
\end{nexample}
\end{examplewrap}

\begin{exercisewrap}
\begin{nexercise}
Verify that the scenario where Brittany is the only one
to exceed the deductible has probability
$(\insureSprob{})^3 (\insureFprob{})^1$.~\footnotemark{}
\end{nexercise}
\end{exercisewrap}
\footnotetext{
  $P(A=\text{\insureS{}},
      \text{ }B=\text{\insureF{}},
      \text{ }C=\text{\insureS{}},
      \text{ }D=\text{\insureS{}})
    = (\insureSprob{})(\insureFprob{})
        (\insureSprob{})(\insureSprob{})
    = (\insureSprob{})^3 (\insureFprob{})^1$.}


The scenario outlined in Example~\ref{insureOneOfFourExceedsDeductible} is an
example of a binomial distribution scenario.
The \termsub{binomial distribution}{distribution!binomial}
describes the probability of having exactly $k$ successes
in $n$ independent Bernoulli trials with probability
of a success $p$
(in Example~\ref{insureOneOfFourExceedsDeductible},
$n=4$, $k=3$, $p=\insureSprob{}$).
We would like to determine the probabilities associated
with the binomial distribution more generally,
i.e. we want a formula where we can use $n$, $k$, and $p$
to obtain the probability.
To do this, we reexamine each part of
Example~\ref{insureOneOfFourExceedsDeductible}.

There were four individuals who could have been the one
to exceed the deductible, and each of these four scenarios
had the same probability.
Thus, we could identify the final probability as
\begin{align*}
[\text{\# of scenarios}] \times P(\text{single scenario})
\end{align*}
The first component of this equation is the number of ways
to arrange the $k=3$ successes among the $n=4$ trials.
The second component is the probability of any of the four
(equally probable) scenarios.

\D{\newpage}

Consider $P($single scenario$)$ under the general case of
$k$ successes and $n-k$ failures in the $n$ trials.
In any such scenario, we apply the Multiplication Rule
for independent events:
\begin{align*}
p^k (1 - p)^{n - k}
\end{align*}
This is our general formula for $P($single scenario$)$.

Secondly, we introduce a general formula for the number
of ways to choose $k$ successes in $n$ trials,
i.e. arrange $k$ successes and $n - k$ failures:
\begin{align*}
{n\choose k} = \frac{n!}{k! (n - k)!}
\end{align*}
The quantity ${n\choose k}$ is read
\term{n choose k}.\footnote{Other notation for
  $n$ choose $k$ includes $_nC_k$, $C_n^k$, and $C(n,k)$.}
The exclamation point notation (e.g. $k!$) denotes
a \term{factorial} expression.\label{factorial_defined}
\begin{align*}
& 0! = 1 \\
& 1! = 1 \\
& 2! = 2\times1 = 2 \\
& 3! = 3\times2\times1 = 6 \\
& 4! = 4\times3\times2\times1 = 24 \\
& \vdots \\
& n! = n\times(n-1)\times...\times3\times2\times1
\end{align*}
Using the formula, we can compute the number of ways
to choose $k = 3$ successes in $n = 4$ trials:
\begin{align*}
{4 \choose 3} = \frac{4!}{3!(4-3)!}
  = \frac{4!}{3!1!} 
  = \frac{4\times3\times2\times1}{(3\times2\times1) (1)}
  = 4
\end{align*}
This result is exactly what we found by carefully thinking
of each possible scenario in
Example~\ref{insureOneOfFourExceedsDeductible}.

Substituting $n$ choose $k$ for the number of scenarios
and $p^k(1-p)^{n-k}$ for the single scenario probability
yields the general binomial formula.

\begin{onebox}{Binomial distribution}
  Suppose the probability of a single trial being
  a success is $p$.
  Then the probability of observing exactly $k$ successes
  in $n$ independent trials is given by\vspace{-1mm}
  \begin{align*}
  {n\choose k}p^k(1-p)^{n-k} = \frac{n!}{k!(n-k)!}p^k(1-p)^{n-k}
  \end{align*}
  The mean, variance, and standard deviation
  of the number of observed successes are\vspace{-2mm}
  \begin{align*}
  \mu &= np
  &\sigma^2 &= np(1-p)
  &\sigma&= \sqrt{np(1-p)}
  \end{align*}
\end{onebox}

\begin{onebox}{Is it binomial? Four conditions to check.}
  \label{isItBinomialTipBox}%
  (1) The trials are independent. \\
  (2) The number of trials, $n$, is fixed. \\
  (3) Each trial outcome can be classified as a \emph{success}
      or \emph{failure}. \\
  (4) The probability of a success, $p$, is the same for
      each trial.
\end{onebox}

\D{\newpage}

\begin{examplewrap}
\begin{nexample}{What is the probability that 3 of 8 randomly
    selected individuals will have exceeded the insurance
    deductible, i.e. that 5 of 8 will not exceed the deductible?
    Recall that 70\% of individuals will not exceed the
    deductible.}
  We would like to apply the binomial model,
  so we check the conditions.
  The number of trials is fixed ($n = 8$) (condition 2)
  and each trial outcome can be classified as a success
  or failure (condition 3).
  Because the sample is random, the trials are independent
  (condition~1) and the probability of a success is the same
  for each trial (condition~4).

  In the outcome of interest, there are $k = 5$ successes
  in $n = 8$ trials (recall that a success is an individual
  who does \emph{not} exceed the deductible), and the
  probability of a success is $p = \insureSprob{}$.
  So the probability that 5 of 8 will not exceed the
  deductible and 3 will exceed the deductible is given by
  \begin{align*}
  { 8 \choose 5}(\insureSprob{})^5
  (1-\insureSprob{})^{8-5}
    &= \frac{8!}{5!(8-5)!}
        (\insureSprob{})^5(1-\insureSprob{})^{8-5} \\
    &= \frac{8!}{5!3!}
        (\insureSprob{})^5(\insureFprob{})^3
  \end{align*}
  Dealing with the factorial part:
  \begin{align*}
  \frac{8!}{5!3!}
    = \frac{8\times7\times6\times5\times4\times3\times2\times1}
        {(5\times4\times3\times2\times1)(3\times2\times1)}
    = \frac{8\times7\times6}{3\times2\times1}
    = 56
  \end{align*}
  Using $(\insureSprob{})^5(\insureFprob{})^3
    \approx \insureBinomEinHSingleScenario{}$,
  the final probability is about
  $56 \times \insureBinomEinHSingleScenario{}
    \approx \insureBinomEinH{}$.
\end{nexample}
\end{examplewrap}

\begin{onebox}{Computing binomial probabilities}
  The first step in using the binomial model is to check
  that the model is appropriate.
  The second step is to identify $n$, $p$, and $k$.
  As the last stage use software or the formulas
  to determine the probability, then interpret the results.%
  \vspace{3mm}

  If you must do calculations by hand, it's often useful
  to cancel out as many terms as possible in the top and
  bottom of the binomial coefficient.
\end{onebox}

\begin{exercisewrap}
\begin{nexercise}
If we randomly sampled 40 case files from the insurance agency
discussed earlier, how many of the cases would you expect to not
have exceeded the deductible in a given year?
What is the standard deviation of the number that would not
have exceeded the deductible?\footnotemark{}
\end{nexercise}
\end{exercisewrap}
\footnotetext{We are asked to determine the expected number
  (the mean) and the standard deviation, both of which can
  be directly computed from the formulas:
  $\mu = np = 40 \times \insureSprob{}
    = \insureBinomFourtyExpValue$
  and $\sigma = \sqrt{np(1-p)}
    = \sqrt{40\times \insureSprob{}\times \insureFprob{}}
    = \insureBinomFourtySD{}$.
  Because very roughly 95\% of observations fall within
  2~standard deviations of the mean
  (see Section~\ref{variability}), we would probably observe
  at least \insureBinomFourtyLower{}
  but fewer than \insureBinomFourtyUpper{} individuals
  in our sample who would not exceed the deductible.}

\begin{exercisewrap}
\begin{nexercise}
The probability that a random smoker will develop a severe
lung condition in his or her lifetime is about $0.3$.
If you have 4 friends who smoke, are the conditions for the
binomial model satisfied?\footnotemark{}
\end{nexercise}
\end{exercisewrap}
\footnotetext{One possible answer:
  if the friends know each other, then the independence
  assumption is probably not satisfied.
  For example, acquaintances may have similar smoking habits,
  or those friends might make a pact to quit together.}

\D{\newpage}

\begin{exercisewrap}
\begin{nexercise}
\label{noMoreThanOneFriendWSevereLungCondition}%
Suppose these four friends do not know each other
and we can treat them as if they were a random sample
from the population.
Is the binomial model appropriate?
What is the probability that\footnotemark{}
\begin{enumerate}[(a)]
\setlength{\itemsep}{0mm}
\item
    None of them will develop a severe lung condition?
\item
    One will develop a severe lung condition?
\item
    That no more than one will develop a severe lung condition?
\end{enumerate}
\end{nexercise}
\end{exercisewrap}
\footnotetext{To check if the binomial model is appropriate,
  we must verify the conditions.
  (i)~Since we are supposing we can treat the friends
  as a random sample, they are independent.
  (ii)~We have a fixed number of trials ($n=4$).
  (iii)~Each outcome is a success or failure.
  (iv)~The probability of a success is the same for each
  trials since the individuals are like a random sample
  ($p=0.3$ if we say a ``success'' is someone getting
  a lung condition, a morbid choice).
  Compute parts~(a) and~(b) using the binomial formula:
  $P(0)
    = {4 \choose 0} (0.3)^0 (0.7)^4
    = 1\times1\times0.7^4
    = 0.2401$,
  $P(1)
    = {4 \choose 1} (0.3)^1(0.7)^{3}
    = 0.4116$.
  Note: $0!=1$.
  Part~(c) can be computed as the sum of parts~(a) and~(b):
  $P(0) + P(1) = 0.2401 + 0.4116 = 0.6517$.
  That is, there is about a 65\% chance that no more than
  one of your four smoking friends will develop a severe
  lung condition.}

\begin{exercisewrap}
\begin{nexercise}
What is the probability that at least 2 of your 4 smoking
friends will develop a severe lung condition in their
lifetimes?\footnotemark{}
\end{nexercise}
\end{exercisewrap}
\footnotetext{The complement (no more than one will develop
  a severe lung condition) as computed in Guided
  Practice~\ref{noMoreThanOneFriendWSevereLungCondition}
  as 0.6517, so we compute one minus this value:~0.3483.}

\begin{exercisewrap}
\begin{nexercise}
Suppose you have 7 friends who are smokers and they can
be treated as a random sample of smokers.\footnotemark{}
\begin{enumerate}[(a)]
\setlength{\itemsep}{0mm}
\item
    How many would you expect to develop a severe lung
    condition, i.e. what is the mean?
\item
    What is the probability that at most 2 of your 7
    friends will develop a severe lung condition.
\end{enumerate}
\end{nexercise}
\end{exercisewrap}
\footnotetext{(a)~$\mu=0.3\times7 = 2.1$.
  (b)~$P($0, 1, or 2 develop severe lung condition$)
      = P(k=0) + P(k=1)+P(k=2) = 0.6471$.}

Next we consider the first term in the binomial probability,
$n$ choose $k$ under some special scenarios.

\begin{exercisewrap}
\begin{nexercise}
Why is it true that ${n \choose 0}=1$ and ${n \choose n}=1$
for any number $n$?\footnotemark{}
\end{nexercise}
\end{exercisewrap}
\footnotetext{Frame these expressions into words.
  How many different ways are there to arrange 0 successes
  and $n$ failures in $n$ trials?
  (1 way.)
  How many different ways are there to arrange $n$ successes
  and 0 failures in $n$ trials?
  (1 way.)}

\begin{exercisewrap}
\begin{nexercise}
How many ways can you arrange one success and $n-1$ failures
in $n$ trials?
How many ways can you arrange $n-1$ successes and one failure
in $n$ trials?\footnotemark{}
\end{nexercise}
\end{exercisewrap}
\footnotetext{One success and $n-1$ failures:
  there are exactly $n$ unique places we can put
  the success, so there are $n$ ways to arrange one
  success and $n-1$ failures.
  A~similar argument is used for the second question.
  Mathematically, we show these results by verifying
  the following two equations:
  \begin{align*}
  {n \choose 1} = n,
    \qquad {n \choose n-1} = n
  \end{align*}}


\newpage


\subsection{Normal approximation to the binomial distribution}
\label{normalApproxBinomialDistSubsection}

\index{distribution!binomial!normal approximation|(}

The binomial formula is cumbersome when the sample size ($n$) is large, particularly when we consider a range of observations. In some cases we may use the normal distribution as an easier and faster way to estimate binomial probabilities.

\newcommand{\smokeprop}{0.15}
\newcommand{\smokeperc}{15\%}
\newcommand{\smokepropcomp}{0.85}
\newcommand{\smokeperccomp}{85\%}
\newcommand{\smokex}{42}
\newcommand{\smokexplusone}{43}
\newcommand{\smoken}{400}
\newcommand{\smokelowertailbinom}{0.0054}
\newcommand{\smokemean}{60}
\newcommand{\smokemeancomp}{340}
\newcommand{\smokesd}{7.14}
\newcommand{\smokez}{-2.52}
\newcommand{\smokelowertailnormal}{0.0059}

\begin{examplewrap}
\begin{nexample}{Approximately \smokeperc{} of the
    US population smokes cigarettes.
    A local government believed their community had
    a lower smoker rate and commissioned a survey of
    400 randomly selected individuals.
    The survey found that only \smokex{} of the
    \smoken{} participants smoke cigarettes.
    If the true proportion of smokers in the community
    was really \smokeperc{}, what is the probability
    of observing \smokex{} or fewer smokers in a sample
    of \smoken{} people?}
  \label{exactBinomSmokerExSetup}%
  We leave the usual verification that the four conditions
  for the binomial model are valid as an exercise.

  The question posed is equivalent to asking,
  what is the probability of observing
  $k=0$, 1, 2, ..., or \smokex{} smokers in a sample of
  $n = \smoken{}$ when $p=\smokeprop{}$?
  We can compute these \smokexplusone{} different
  probabilities and add them together to find the answer:
  \begin{align*}
  &P(k=0\text{ or }k=1\text{ or }\cdots\text{ or } k=\smokex{}) \\
	&\qquad = P(k=0) + P(k=1) + \cdots + P(k=\smokex{}) \\
	&\qquad = \smokelowertailbinom{}
  \end{align*}
  If the true proportion of smokers in the community
  is $p=\smokeprop{}$, then the probability of observing
  \smokex{} or fewer smokers in a sample of $n=\smoken{}$
  is \smokelowertailbinom{}.
\end{nexample}
\end{examplewrap}

The computations in Example~\ref{exactBinomSmokerExSetup}
are tedious and long.
In general, we should avoid such work if an alternative method
exists that is faster, easier, and still accurate.
Recall that calculating probabilities of a range of values
is much easier in the normal model.
We might wonder, is it reasonable to use the normal model
in place of the binomial distribution?
Surprisingly, yes, if certain conditions are met.

\begin{exercisewrap}
\begin{nexercise}
Here we consider the binomial model when the probability
of a success is $p = 0.10$.
Figure~\ref{fourBinomialModelsShowingApproxToNormal}
shows four hollow histograms for simulated samples from
the binomial distribution using four different sample sizes:
$n = 10, 30, 100, 300$.
What happens to the shape of the distributions as the sample
size increases?
What distribution does the last hollow histogram
resemble?\footnotemark{}
\end{nexercise}
\end{exercisewrap}
\footnotetext{The distribution is transformed from a blocky
  and skewed distribution into one that rather resembles the
  normal distribution in last hollow histogram.}

\begin{figure}[h]
  \centering
  \Figure[Four hollow histograms are shown, each in their own plot, based on a probability of p equals 0.10 and sample sizes of n equals 10, 30, 100, and 300. The first plot for n = 10 shows a distribution centered at 1 and is notably right skewed. The second plot for n = 30 shows a distribution centered at about 3, is just a bit right skewed, and is starting to look a little bit like a bell-shaped distribution. The third plot for n = 100 shows a distribution centered at about 10 and that is almost entirely symmetric with just the slightest indication it is right skewed. This third distribution also looks very bell-shaped. The fourth plot for n = 300 shows a distribution centered at about 30 and that is symmetric. This last plot looks very bell-shaped and resembles a normal distribution.]{0.92}{fourBinomialModelsShowingApproxToNormal}
  \caption{Hollow histograms of samples from the binomial
      model when $p = 0.10$.
      The sample sizes for the four plots are
      $n = 10$, 30, 100, and 300, respectively.}
  \label{fourBinomialModelsShowingApproxToNormal}
\end{figure}

\begin{onebox}{Normal approximation of the binomial distribution}
  The binomial distribution with probability of success
  $p$ is nearly normal when the sample size $n$ is
  sufficiently large that $np$ and $n(1-p)$ are both
  at least 10.
  The approximate normal distribution has parameters
  corresponding to the mean and standard deviation of
  the binomial distribution:\vspace{-1.5mm}
  \begin{align*}
  \mu &= np
      &\sigma& = \sqrt{np(1 - p)}
  \end{align*}
\end{onebox}

The normal approximation may be used when computing
the range of many possible successes.
For instance, we may apply the normal distribution to
the setting of Example~\ref{exactBinomSmokerExSetup}.

\D{\newpage}

\begin{examplewrap}
\begin{nexample}{How can we use the normal approximation
    to estimate the probability of observing \smokex{} or
    fewer smokers in a sample of \smoken{}, if the true
    proportion of smokers is $p = \smokeprop{}$?}
  \label{approxNormalForSmokerBinomEx}
  Showing that the binomial model is reasonable was a
  suggested exercise in Example~\ref{exactBinomSmokerExSetup}.
  We also verify that both $np$ and $n(1-p)$ are at least 10:
  \begin{align*}
  np &= \smoken{} \times \smokeprop{} = \smokemean{}
  &n (1 - p) &= \smoken{} \times \smokepropcomp{}
      = \smokemeancomp{}
  \end{align*}
  With these conditions checked, we may use the normal
  approximation in place of the binomial distribution
  using the mean and standard deviation from the binomial
  model:
  \begin{align*}
  \mu &= np = \smokemean{}
  &\sigma &= \sqrt{np(1 - p)} = \smokesd{}
  \end{align*}
  We want to find the probability of observing
  \smokex{} or fewer smokers using this model.
\end{nexample}
\end{examplewrap}

\begin{exercisewrap}
\begin{nexercise}
Use the normal model $N(\mu = \smokemean{}, \sigma = \smokesd{})$
to estimate the probability of observing \smokex{} or fewer
smokers.
Your answer should be approximately equal to the solution
of Example~\ref{exactBinomSmokerExSetup}:%
~\smokelowertailbinom{}.~\footnotemark{}
\end{nexercise}
\end{exercisewrap}
\footnotetext{Compute the Z-score first:
  $Z = \frac{\smokex{} - \smokemean{}}{\smokesd{}} = \smokez{}$.
  The corresponding left tail area is \smokelowertailnormal{}.}



\newpage


\subsection{The normal approximation breaks down on small intervals}

The normal approximation to the binomial distribution tends to perform poorly when estimating the probability of a small range of counts, even when the conditions are met.

\newcommand{\smokeA}{49}
\newcommand{\smokeB}{50}
\newcommand{\smokeC}{51}
\newcommand{\smokeABCBinom}{0.0649}
\newcommand{\smokeABCNormal}{0.0421}
\newcommand{\smokeABCNormalFixed}{0.0633}

Suppose we wanted to compute the probability of observing
\smokeA{}, \smokeB{}, or \smokeC{} smokers in \smoken{}
when $p = \smokeprop{}$.
With such a large sample, we might be tempted to apply
the normal approximation and use the range \smokeA{} to \smokeC{}.
However, we would find that the binomial solution and the normal
approximation notably differ:
\begin{align*}
\text{Binomial:}&\ \smokeABCBinom{}
&\text{Normal:}&\ \smokeABCNormal{}
\end{align*}
We can identify the cause of this discrepancy using
Figure~\ref{normApproxToBinomFail}, which shows the areas
representing the binomial probability (outlined) and normal
approximation (shaded).
Notice that the width of the area under the normal
distribution is 0.5 units too slim on both sides of
the interval.

\begin{figure}[h]
  \centering
  \Figure[A normal distribution centered at 60 with a standard deviation of about 7 is shown. (The determination that the standard deviation is about 7 was based on the normal distribution being very close to 0 a distance of about 20 from the mean, and this happens about 3 standard deviations from the mean.) A region of this distribution is shaded from 49 to 51. Additionally, a red outlined area is boxed out between 48.5 and 51.5 that represents the exact binomial distribution.]{1.0}{normApproxToBinomFail}
  \caption{A normal curve with the area between
      \smokeA{} and \smokeC{} shaded.
      The outlined area represents the exact binomial
      probability.}
  \label{normApproxToBinomFail}
\end{figure}

\begin{onebox}{Improving the normal approximation
    for the binomial distribution}
  The normal approximation to the binomial distribution
  for intervals of values is usually improved if cutoff
  values are modified slightly.
  The cutoff values for the lower end of a shaded region
  should be reduced by 0.5, and the cutoff value for the
  upper end should be increased by 0.5.
\end{onebox}

The tip to add extra area when applying the normal
approximation is most often useful when examining
a range of observations.
In the example above, the revised normal distribution
estimate is \smokeABCNormalFixed{}, much closer to the
exact value of \smokeABCBinom{}.
While it is possible to also apply this correction when
computing a tail area, the benefit of the modification
usually disappears since the total interval is typically
quite wide.

\index{distribution!binomial!normal approximation|)}
\index{distribution!binomial|)}


{\input{ch_distributions/TeX/binomial_distribution.tex}}




%_________________
